% SOURCE_STATEMENT: "For the periodic Lorentz gas with scatterer radius $r$ and time $t$, can one establish a joint limit as $r\\to0$ and $t\\to\\infty$, rather than taking the limits successively?"
% Preparation and final self-audit: CONTRIBUTING.md prompts C and D; source/audit artifacts retained in the local physics_statements_2026_10_11 batch.
% !TeX program = lualatex
\documentclass[11pt]{article}
\usepackage[margin=25mm]{geometry}
\usepackage{fontspec}
\setmainfont{Latin Modern Roman}
\usepackage{amsmath,amssymb,mathtools}
\usepackage{unicode-math}
\setmathfont{Latin Modern Math}
\usepackage{xurl,hyperref}
\hypersetup{colorlinks=true,urlcolor=blue}
\setcounter{secnumdepth}{0}
\setlength{\parindent}{0pt}
\setlength{\parskip}{5pt}
\setlength{\emergencystretch}{3em}
\newcommand{\sref}[2]{\hyperlink{source:#1:#2}{[S#2]}}
\newcommand{\cataloguescope}[1]{\paragraph{Recorded target.}#1}
\newcommand{\R}{\mathbb R}
\newcommand{\C}{\mathbb C}
\newcommand{\N}{\mathbb N}
\newcommand{\Z}{\mathbb Z}
\newcommand{\Q}{\mathbb Q}
\newcommand{\F}{\mathbb F}
\newcommand{\eps}{\varepsilon}
\begin{document}
\section{Joint Boltzmann–Grad and Long-Time Limits for Lorentz Gases}\label{statement}

\subsection{Definitions and mathematical statement}

\paragraph{Shared notation.}$\mathbb N=\{1,2,\ldots\}$; the standard number fields have their usual meanings. All additional conventions are stated locally.


\paragraph{Definitions and conventions.}
A periodic Lorentz gas is a particle moving freely between elastic collisions with periodically arranged fixed scatterers. The Boltzmann--Grad regime shrinks scatterer radius $r$ while rescaling path length or time to keep collision statistics nontrivial. A long-time limit studies transport after many collisions. A joint limit controls these approximations as $r\to0$ and $t\to\infty$ together, rather than applying one asymptotic result after taking the other limit.

\paragraph{Target retained from the canonical source catalogue.}
\paragraph{Catalogue wording retained for traceability.}
For the periodic Lorentz gas with scatterer radius $r$ and time $t$, can one establish a joint limit as $r\to0$ and $t\to\infty$, rather than taking the limits successively?

\paragraph{Source scope and qualification.}
The joint scaling relation, normalization and limiting observable are unspecified in the extraction. They must come from the report; unrestricted convergence along every possible pair $(r,t)$ is not silently asserted.
\subsection{Short English statement}
Can dilute-scatterer and long-time asymptotics for the periodic Lorentz gas be justified simultaneously?
\subsection{Sources}
\begin{itemize}
\item[\hypertarget{source:statement:1}{S1}] ULAM.AI and the UnsolvedMath Contributors, UnsolvedMath, catalogue entry. Catalogue curation under CC BY 4.0; definitions and English exposition added locally.
\url{https://huggingface.co/datasets/ulamai/UnsolvedMath}.
\item[\hypertarget{source:statement:2}{S2}] Scaling Limits in Models of Statistical Mechanics (2015), Oberwolfach Reports, DOI 10.4171/owr/2015/40, pp. 2329--2398. Primary report accessed 2026-10-11; relevant PDF pages 32.
\url{https://doi.org/10.4171/owr/2015/40}.
\item[\hypertarget{source:statement:3}{S3}] Marklof--Toth, Superdiffusion in the Periodic Lorentz Gas (catalogue-listed context; not independently consulted in this preparation).
\url{https://doi.org/10.1007/s00220-016-2578-y}.
\item[\hypertarget{source:statement:4}{S4}] Boca--Zaharescu, Periodic Lorentz Gas with Small Scatterers (catalogue-listed context; not independently consulted in this preparation).
\url{https://doi.org/10.1007/s00440-023-01197-6}.
\end{itemize}
\label{end:statement}
\end{document}
